\section{Methodology}
\label{sec:methodology}

The proposed solution implements a comprehensive dual-subsystem approach integrating mechanical and electrical innovations to address identified system vulnerabilities. This chapter presents the design methodology and implementation strategy for both subsystems.

\subsection{System Architecture Overview}

The N4 recovery system integrates five primary subsystems working in coordinated fashion to achieve reliable dual deployment:

\begin{enumerate}
    \item \textbf{Telemetry System:} Long-range data transmission using XBee-PRO 900HP modules
    \item \textbf{Ignition Control System:} PWM-controlled nichrome wire deployment igniters
    \item \textbf{Feedback Monitoring System:} Real-time battery and igniter continuity sensing
    \item \textbf{Mechanical Deployment System:} Sealed explosion cap with pyrotechnic charge containment
    \item \textbf{Avionics Support Structure:} Ruggedized 3D-printed component housing
\end{enumerate}

Figure~\ref{fig:system_architecture} presents the high-level system architecture showing data flow and control signal paths between subsystems.

% Note: Add figure when diagram is created
% \begin{figure}[H]
% \centering
% \includegraphics[width=0.85\textwidth]{system_architecture.pdf}
% \caption{N4 Recovery System Architecture}
% \label{fig:system_architecture}
% \end{figure}

\subsection{Mechanical System Design}

\subsubsection{Deployment Mechanism: The Sealed Explosion Cap}

The sealed explosion cap design addresses two critical requirements: containment of combustion flame products and consistent force delivery to shear pin interfaces.

\textbf{Design Requirements:}
\begin{itemize}
    \item Withstand internal pressure of 5--10 bar generated by black powder combustion
    \item Prevent direct flame contact with parachute materials
    \item Direct pressure wave energy to shear pin interfaces
    \item Maintain structural integrity under launch acceleration
    \item Enable repeatable assembly/disassembly for testing
\end{itemize}

\textbf{Material Selection:} The explosion cap will be fabricated from aluminum 6061-T6, selected for its combination of strength (yield strength $\sim$275 MPa), machinability, and thermal conductivity. Wall thickness is designed according to thin-wall pressure vessel theory:

\begin{equation}
t = \frac{PD}{2\sigma_{allow}}
\end{equation}

Where $t$ is wall thickness, $P$ is internal pressure, $D$ is diameter, and $\sigma_{allow}$ is allowable stress (applying appropriate safety factor to yield strength).

\textbf{Sealing Mechanism:} O-ring sealing using Buna-N (nitrile) rubber provides gas-tight containment while enabling repeated disassembly. O-ring groove dimensions follow AS568A standard specifications to ensure proper compression and sealing performance.

\subsubsection{Charge Sizing and Ejection Force Characterization}

One critical research gap addresses in this project is establishing quantifiable relationships between charge mass and ejection force for the N4 airframe geometry.

\textbf{Experimental Approach:}
\begin{enumerate}
    \item Static force testing using load cell measurements
    \item Varying charge masses from 1--5 grams in 0.5 gram increments
    \item Recording peak force, rise time, and pressure-time profiles
    \item Characterizing charge-to-force relationship through regression analysis
\end{enumerate}

Expected outcome: Empirical data enabling precise charge sizing to achieve target ejection force of 200--300 N, sufficient to shear deployment pins without over-stressing airframe structure.

\subsubsection{Avionics Bay Structure}

Design and fabrication of a custom 3D-printed sled/holder assembly to secure critical flight computer components, LiPo battery packs, and Battery Management System (BMS) modules.

\textbf{Design Requirements:}
\begin{itemize}
    \item Secure mounting for ESP32 flight computer, XBee module, GPS, and BMS
    \item Vibration isolation to reduce high-frequency transmission to PCB
    \item Maintain center of gravity within 5 mm of airframe centerline
    \item Withstand 15G sustained acceleration and 100G shock loads
    \item Facilitate cable routing and component access
\end{itemize}

\textbf{Vibration Isolation Strategy:} The holder incorporates compliant flexure elements designed to provide mechanical decoupling at frequencies above 200 Hz (primary motor vibration frequency). Finite element analysis (FEA) will validate the design under expected loading conditions.

\textbf{Material Selection:} PETG (Polyethylene Terephthalate Glycol-modified) is selected for flight hardware fabrication due to:
\begin{itemize}
    \item Superior layer adhesion compared to PLA
    \item Operating temperature range: -40°C to +70°C
    \item Impact resistance 5--7× greater than PLA
    \item Low moisture absorption
\end{itemize}

\subsection{Electrical System Architecture}

\subsubsection{Hardware Architecture}

The electrical system architecture comprises four primary functional blocks:

\begin{enumerate}
    \item \textbf{Microcontroller Unit (MCU):} ESP32 platform
    \item \textbf{Telemetry Module:} XBee-PRO 900HP
    \item \textbf{Sensor Suite:} GPS, accelerometer, barometer
    \item \textbf{Power Management:} BMS and voltage monitoring
    \item \textbf{Actuation System:} MOSFET-driven ignition circuits
\end{enumerate}

\textbf{ESP32 Selection Rationale:}
\begin{itemize}
    \item Dual-core Xtensa LX6 processor (240 MHz)
    \item 520 KB SRAM, 4 MB Flash
    \item Hardware UART, SPI, I$^2$C interfaces
    \item Built-in WiFi/Bluetooth (for ground testing and configuration)
    \item Low power consumption with sleep modes
    \item Extensive third-party library support
\end{itemize}

\subsubsection{PCB Design}

Custom printed circuit board design addresses several failure modes observed in previous implementations:

\textbf{Design Objectives:}
\begin{itemize}
    \item Eliminate loose wire connections susceptible to vibration failure
    \item Minimize circuit footprint to fit within 70 × 70 mm avionics bay envelope
    \item Implement robust power distribution with appropriate decoupling
    \item Provide test points for debugging and validation
    \item Support modular replacement of key components (ESP32, XBee)
\end{itemize}

\textbf{PCB Specifications:}
\begin{itemize}
    \item 2-layer FR-4 substrate (1.6 mm thickness)
    \item 2 oz copper weight on power traces
    \item Minimum trace width: 0.3 mm (signal), 1.0 mm (power)
    \item ENIG (Electroless Nickel Immersion Gold) surface finish for reliability
\end{itemize}

% Note: Add figure when PCB design is complete
% \begin{figure}[H]
% \centering
% \includegraphics[width=0.85\textwidth]{pcb_layout.pdf}
% \caption{Custom PCB Layout for N4 Recovery System}
% \label{fig:pcb_layout}
% \end{figure}

\subsubsection{Smart Ignition Control System}

Implementation of Pulse Width Modulation (PWM) controlled ignition utilizing low-side MOSFET driver architecture addresses the thermal runaway problem identified in the literature review.

\textbf{Circuit Topology:}
The ignition circuit employs an IRLZ44N logic-level N-channel MOSFET driven directly from ESP32 PWM output. Key specifications:
\begin{itemize}
    \item $V_{DS(max)}$: 55V
    \item $I_D$: 47A continuous
    \item $R_{DS(on)}$: 22 m$\Omega$ @ $V_{GS}$ = 5V
    \item Logic-level gate threshold: $V_{GS(th)}$ = 1--2V
\end{itemize}

\textbf{PWM Control Algorithm:}
The system implements a closed-loop temperature control algorithm:
\begin{enumerate}
    \item Initialize PWM at low duty cycle (10--20\%)
    \item Ramp duty cycle linearly over 2--5 second interval
    \item Monitor wire resistance (proportional to temperature) via voltage sensing
    \item Hold at optimal ignition temperature (800--900°C) until black powder ignites
    \item Detect ignition via accelerometer spike and terminate PWM output
\end{enumerate}

This approach prevents wire burnout while ensuring reliable ignition, enabling multiple ground tests with the same nichrome wire element.

\textbf{Protection Features:}
\begin{itemize}
    \item Over-current detection via current sense resistor
    \item Thermal shutdown if nichrome resistance exceeds failure threshold
    \item Watchdog timer preventing sustained ignition (>10 seconds)
    \item Manual enable requirement (dual-switch safety)
\end{itemize}

\subsubsection{Power Management and Feedback Systems}

\textbf{Battery Management System Integration:}
The system employs a dedicated 4S LiPo Battery Management System (BMS) providing:
\begin{itemize}
    \item Cell balancing during charging
    \item Over-voltage protection (4.25V per cell)
    \item Under-voltage protection (3.0V per cell)
    \item Over-current protection (20A continuous limit)
    \item Short-circuit protection
\end{itemize}

\textbf{Voltage Monitoring:}
ESP32 ADC channels monitor battery voltage through resistive divider networks. The voltage divider is designed to scale the 12--16.8V battery range to the ESP32's 0--3.3V ADC input range. Voltage is sampled at 10 Hz and transmitted to ground station in real-time.

\textbf{Igniter Continuity Detection:}
Continuity of nichrome wire igniters is verified by passing a small test current (10--50 mA) and measuring voltage drop. The system reports:
\begin{itemize}
    \item \textbf{GOOD:} Resistance in range 2--10 $\Omega$
    \item \textbf{OPEN:} Resistance >100 $\Omega$ (wire break or disconnection)
    \item \textbf{SHORT:} Resistance <0.5 $\Omega$ (degraded connection)
\end{itemize}

\textbf{Launch Inhibit Safety Protocol:}
The system implements automatic "Launch Inhibit" safety protocol that prevents launch when:
\begin{itemize}
    \item Main battery voltage <14.0V (insufficient for reliable deployment)
    \item Any igniter reports OPEN or SHORT status
    \item XBee telemetry link not established
    \item GPS fix not acquired (optional, configurable)
\end{itemize}

Visual and audible indicators on the ground station alert operators to inhibit conditions.

\subsection{Software Architecture}

\subsubsection{Flight Computer Firmware}

The ESP32 firmware implements a task-based architecture using FreeRTOS:

\textbf{Primary Tasks:}
\begin{enumerate}
    \item \textbf{Sensor Task:} Read GPS, accelerometer, barometer at 10 Hz
    \item \textbf{Telemetry Task:} Package and transmit data via XBee at 5 Hz
    \item \textbf{State Machine Task:} Manage flight state (Pre-launch, Powered Ascent, Coasting, Apogee, Descent, Landed)
    \item \textbf{Deployment Task:} Monitor altitude/acceleration for deployment triggers
    \item \textbf{Ignition Control Task:} Execute PWM control algorithm for deployment actuation
\end{enumerate}

\subsection{Testing Methodology}

Comprehensive testing validates each subsystem independently before system integration:

\subsubsection{Component-Level Testing}
\begin{itemize}
    \item PWM ignition circuit: Verify duty cycle accuracy, current delivery, thermal response
    \item Telemetry system: Range testing, packet loss characterization, latency measurement
    \item Voltage monitoring: Accuracy verification across battery voltage range
    \item Continuity detection: Resistance measurement accuracy, fault detection
\end{itemize}

\subsubsection{Subsystem Integration Testing}
\begin{itemize}
    \item Static ignition tests: Validate PWM control prevents wire burnout
    \item Deployment force testing: Characterize charge-to-force relationship
    \item Vibration table testing: Validate avionics holder isolation performance
    \item Drop testing: Simulate deployment shock loads
\end{itemize}

\subsubsection{System-Level Testing}
\begin{itemize}
    \item Complete recovery system ground tests
    \item Telemetry range validation at test field
    \item Simulated flight profile testing
    \item Final flight readiness review
\end{itemize}

\subsection{Summary}

This chapter presented the comprehensive methodology for design, fabrication, and testing of the N4 dual recovery system. The mechanical system addresses deployment reliability through sealed explosion cap design and ruggedized avionics structure. The electrical system implements XBee-PRO telemetry, PWM-controlled ignition, and comprehensive feedback monitoring. Systematic testing validates each subsystem before integration. The next chapter presents results obtained from initial design and testing activities.
